\section{System Overview}
The face recognition system follows a modular architecture design pattern, separating concerns into distinct components that can work together seamlessly while allowing for independent development, testing, and replacement. Figure \ref{fig:architecture} illustrates the high-level architecture of the system.

\begin{figure}[H]
\centering
\includegraphics[width=0.9\textwidth]{images/system_architecture.png}
\caption{High-level system architecture showing the main components and their interactions}
\label{fig:architecture}
\end{figure}

The system is composed of several key modules:

\begin{itemize}
    \item \textbf{Core Components:} Face detection, recognition, verification, alignment, and drowsiness detection modules
    \item \textbf{Services:} Camera service, image processing service
    \item \textbf{Database:} Storage and retrieval of facial embeddings
    \item \textbf{API:} Web interface for interaction with the system
    \item \textbf{Configuration:} Centralized settings management
    \item \textbf{Utilities:} Shared helper functions for image processing and camera handling
\end{itemize}

\section{Component Interactions}
The system components interact in a pipeline fashion for processing images and video streams:

\begin{enumerate}
    \item The camera handler captures frames from a camera source
    \item Frames are processed by the face detection module to locate faces
    \item Detected faces are aligned using the alignment module
    \item Facial landmarks are extracted for additional processing (e.g., drowsiness detection)
    \item The face recognition module generates embeddings for each detected face
    \item The verification module compares embeddings against the database
    \item Results are displayed through the GUI or returned via the API
\end{enumerate}

\section{Operational Modes}
The system supports three primary operational modes:

\subsection{Single Image Processing}
In this mode, the system processes a single image file, detecting faces, generating embeddings, and comparing against the database. This mode is useful for testing and for processing static images.

\subsection{Camera Feed Processing}
This mode processes a live camera feed in real-time, continuously detecting and recognizing faces. It can display the processed feed with detection results and recognized user names. This mode is suitable for interactive applications like access control systems.

\subsection{API Server Mode}
The API mode exposes the system's capabilities through a web interface, allowing other applications to interact with the face recognition functionality. It provides endpoints for uploading images, retrieving detection results, and managing the user database.

\section{Modularity and Extensibility}
The system's modular design enables several key advantages:

\begin{itemize}
    \item \textbf{Interchangeable Components:} Different detection models (MediaPipe, YOLOv8, etc.) can be swapped without affecting other components
    \item \textbf{Extensibility:} New models and algorithms can be integrated by implementing the defined interfaces
    \item \textbf{Testability:} Components can be tested independently
    \item \textbf{Maintainability:} Changes to one component don't require modifications to others
\end{itemize}

\section{Data Flow}
Figure \ref{fig:dataflow} illustrates the flow of data through the system components during typical operation.

\begin{figure}[H]
\centering
\includegraphics[width=0.8\textwidth]{images/data_flow.png}
\caption{Data flow diagram showing how information moves through the system}
\label{fig:dataflow}
\end{figure}

The system's architecture emphasizes clean separation of concerns, making it easy to understand, maintain, and extend. Each component focuses on a specific task while communicating with other components through well-defined interfaces.